We say that a square-free positive integer \(n\) is \emph{almost prime} if
\[
n \mid x^{d_1} + x^{d_2} + \ldots + x^{d_k} - kx
\]
for all integers \(x\), where \(1 = d_1 < d_2 < \ldots < d_k = n\) are all the positive divisors of \(n\). Suppose that \(r\) is a Fermat prime (i.e., it is a prime of the form \(2^{2^m} + 1\) for an integer \(m \ge 0\)), \(p\) is a prime divisor of an almost prime integer \(n\), and \(p \equiv 1 \pmod{r}\). Show that, with the above notation, \(d_i \equiv 1 \pmod{r}\) for all \(1 \le i \le k\).

(An integer \(n\) is called \emph{square-free} if it is not divisible by \(d^2\) for any integer \(d > 1\).)
